\scanpage{559}
\mspage{\unsure{583}}

\deleted{We set}
\[
\deleted{
\left\{
\begin{aligned}
S\unsure{Z_\rho|N_\rho}
 &=\operatorname{Ker}\bigl(\unsure{Z_\rho|N_\rho}
       \xrightarrow{\delta_\rho}\gamma_{\rho,\sigma}\bigr)
       =\operatorname{Norm}_{\mathcal G}(L_\rho)\\
SN_\sigma
 &=\operatorname{Ker}\bigl(N_\sigma
       \xrightarrow{\delta_\sigma}\gamma_{\rho,\sigma}\bigr)
       =\operatorname{Norm}_{\mathcal G}(L_\sigma)
\end{aligned}
\right.
}
\tag{15}
\]
\[
\deleted{
S\unsure{B}_{\rho,\sigma}
=\operatorname{Ker}\bigl(\unsure{B}_{\rho,\sigma}
   \xrightarrow{\delta_B}\gamma_{\rho,\sigma}\bigr)
\simeq L_0\subset
\left\{
\begin{gathered}
\unsure{\operatorname{Norm}|\operatorname{Cent}}_{\mathcal G}
       (\unsure{L_0|L_\rho})\\[-2pt]
\cap\\[-2pt]
\operatorname{Norm}_{\mathcal G}\,\unsure{L_0|L_\sigma}
\end{gathered}
\right.
}
\]
\margin{Canceled\par
\deleted{\unsure{It seems that}
\unsure{in the case that} \unsure{we want to consider},
\(\operatorname{Norm}_{\mathcal G}(L_0)
\unsure{\ne|=}\unsure{\operatorname{Cent}}_{\mathcal G}(L_0)\)
and idem for \(L_\rho,L_\sigma\),}}

\textcircled{\(5^\circ\)}
\margin{The groups
\(S_0\mathcal M_{\rho,\sigma},\mathcal M_{\rho,\sigma},
S_0\Gamma_{\rho,\sigma},\Gamma_{\rho,\sigma}\).}

We set
\[
\left\{
\begin{aligned}
S_0\mathcal M_{\rho,\sigma}
 &=\deleted{\unsure{B}}Z_{\rho,\sigma}
   \mathop{\times}_{\gamma_{\rho,\sigma}}N_\rho^*
   \mathop{\times}_{\gamma_{\rho,\sigma}}N_\sigma^*
   \mathop{\times}_{\gamma_{\rho,\sigma}}G_{\rho,\sigma}.\\
\mathcal M_{\rho,\sigma}
 &=N_\rho^*
   \mathop{\times}_{\gamma_{\rho,\sigma}}N_\sigma^*
   \mathop{\times}_{\gamma_{\rho,\sigma}}G_{\rho,\sigma}
\end{aligned}
\right.
\tag{50}
\]
\[
\deleted{\unsure{42}}
\qquad
\left\{
\begin{aligned}
S_0\Gamma_{\rho,\sigma}
 &=Z_{\rho,\sigma}
   \mathop{\times}_{\gamma_{\rho,\sigma}}N_\rho^*
   \mathop{\times}_{\gamma_{\rho,\sigma}}N_\sigma^*\\
\Gamma_{\rho,\sigma}
 &=\unsure{N_\rho^*}
   \mathop{\times}_{\gamma_{\rho,\sigma}}N_\sigma^*,
\end{aligned}
\right.
\tag{51}
\]

whence \deleted{\unsure{homomorphism}}\inserted{\uline{diagram}} of exact sequences
\diagram{
\((52)\qquad\)
\begin{tikzcd}[column sep=large,row sep=large]
 & & 1 \arrow[d] & 1 \arrow[d] & \\
 & & L_0 \arrow[d] \arrow[r,equal] & L_0 \arrow[d] & \\
1 \arrow[r] & \mathcal G^+ \arrow[r] \arrow[d,equal]
 & S_0\mathcal M_{\rho,\sigma} \arrow[r] \arrow[d]
 & S_0\Gamma_{\rho,\sigma} \arrow[r] \arrow[d] & 1\\
1 \arrow[r] & \mathcal G^+ \arrow[r]
 & \mathcal M_{\rho,\sigma} \arrow[r] \arrow[d]
 & \Gamma_{\rho,\sigma} \arrow[r] \arrow[d] & 1,\\
 & & 1 & 1 &
\end{tikzcd}
}

and exact sequences
\diagram{
\((53)\qquad\)
\begin{tikzcd}[column sep=large]
1 \arrow[r]
 & \underset{\simeq L_0}{\uline{SZ_{\rho,\sigma}}}
   \times SN_\rho^*\times SN_\sigma^*\times\mathcal G^+
       \arrow[r]
 & S_0\mathcal M_{\rho,\sigma}\arrow[r]
 & \gamma_{\rho,\sigma}\arrow[r] & 1
\end{tikzcd}
}
\diagram{
\((54)\qquad\)
\begin{tikzcd}[column sep=large]
1 \arrow[r]
 & SN_\rho^*\times SN_\sigma^*\times\mathcal G^+\arrow[r]
 & \mathcal M_{\rho,\sigma}\arrow[r]
 & \gamma_{\rho,\sigma}\arrow[r] & 1
\end{tikzcd}
}
\diagram{
\((55)\qquad\)
\begin{tikzcd}[column sep=large]
1 \arrow[r]
 & \underset{\simeq L_0}{\uline{SZ_{\rho,\sigma}}}
   \times SN_\rho^*\times SN_\sigma^*\arrow[r]
 & S_0\Gamma_{\rho,\sigma}\arrow[r]
 & \gamma_{\rho,\sigma}\arrow[r] & 1
\end{tikzcd}
}
\diagram{
\((56)\qquad\)
\begin{tikzcd}[column sep=large]
1 \arrow[r]
 & SN_\rho^*\times SN_\sigma^*\arrow[r]
 & \Gamma_{\rho,\sigma}\arrow[r]
 & \gamma_{\rho,\sigma}\arrow[r] & 1.
\end{tikzcd}
}

\mspage{558}

{\small
\begin{center}
558\\
\texttt{https://grothendieck.umontpellier.fr}\\
Université de Montpellier, UMR CNRS 5149, Case courrier 051, Place Eugène Bataillon\\
34095 Montpellier cedex 5 - France
\end{center}
}

\scanpage{560}

We also have
\[
S_0\mathcal M_{\rho,\sigma}
\simeq\mathcal M_{\rho,\sigma}
       \mathop{\times}_{\gamma_{\rho,\sigma}}Z_{\rho,\sigma}
\tag{57}
\]
\[
S_0\Gamma_{\rho,\sigma}
\simeq\Gamma_{\rho,\sigma}
       \mathop{\times}_{\gamma_{\rho,\sigma}}Z_{\rho,\sigma}
\tag{58}
\]

i.e. the extensions \(S_0\mathcal M_{\rho,\sigma},S_0\Gamma_{\rho,\sigma}\)
of \(\mathcal M_{\rho,\sigma},\Gamma_{\rho,\sigma}\) by \(L_0\)
are deduced from the extension \(Z_{\rho,\sigma}\) of
\(\gamma_{\rho,\sigma}\) by \(L_0\) by inverse image under
\(\mathcal M_{\rho,\sigma}\to\gamma_{\rho,\sigma},
\Gamma_{\rho,\sigma}\to\gamma_{\rho,\sigma}\) --- and we also have
\deleted{\unsure{$\Gamma$}}
\[
S_0\mathcal M_{\rho,\sigma}
\simeq\mathcal M_{\rho,\sigma}
       \mathop{\times}_{\Gamma_{\rho,\sigma}}S_0\Gamma_{\rho,\sigma}
\tag{59}
\]

i.e. the extension \(S_0\mathcal M_{\rho,\sigma}\) of
\(\mathcal M_{\rho,\sigma}\) by \(L_0\) is deduced from the extension
\(S_0\Gamma_{\rho,\sigma}\) of \(\Gamma_{\rho,\sigma}\) by \(L_0\),
by inverse image \unsure{under}
\(\mathcal M_{\rho,\sigma}\to\Gamma_{\rho,\sigma}\).

\inserted{We introduce an element}
\[
\deleted{\unsure{\tau'_{\mathcal M}}},
\qquad
\tau'_{\deleted{\unsure{S_0}}\mathcal M}
  \in\deleted{S_0}\mathcal M_{\rho,\sigma}
\quad
\deleted{\text{\unsure{the}}}
\quad
\tau'_{\deleted{\unsure{S_0}}\mathcal M}
 =\deleted{\unsure{\tau'_Z}}
 \bigl(
  \overset{\inserted{\substack{a\\\Vert}}}{\tau'_\rho},
  \overset{\inserted{\substack{b\\\Vert}}}{\tau'_\sigma},
  \overset{\inserted{\substack{u\\\Vert}}}{\tau'_G}
 \bigr)
\tag{60}
\]

\deleted{\unsure{of homomorphisms} between the groups
introduced, I am going to link them in
\deleted{in}\inserted{a} commutative diagram}
\inserted{where \(\deleted{\unsure{\tau'_Z}},\tau'_\rho,\tau'_\sigma,
\tau'_G\) are the}
``avatars'' of \(\tau'\), taken resp. in
\(\deleted{\unsure{Z_{\rho,\sigma}}}
N_\rho^*,N_\sigma^*,G_{\rho,\sigma}\).
These do indeed all have the same image in
\(G_{\rho,\sigma}\) and a fortiori in \(\gamma_{\rho,\sigma}\),
hence define an \unsure{element} of
\(\deleted{\unsure{G}}\inserted{\mathcal M}_{\rho,\sigma}\),
whose operation on \(\mathcal G^+\) is that of \(\tau\).
Whence an embedding
\diagram{
\((61)\qquad\)
\begin{tikzcd}
\mathcal G \arrow[r,hook,"i_{\mathcal G,\mathcal M}"]
 & \mathcal M_{\rho,\sigma}
\end{tikzcd}
\quad : identity on \(\mathcal G^+\)
\qquad
\begin{tikzcd}
\tau'_{\mathcal G}\arrow[r,mapsto] & \tau'_{\mathcal M}.
\end{tikzcd}
}

\mspage{559}

{\small
\begin{center}
559\\
\texttt{https://grothendieck.umontpellier.fr}\\
Université de Montpellier, UMR CNRS 5149, Case courrier 051, Place Eugène Bataillon\\
34095 Montpellier cedex 5 - France
\end{center}
}

\scanpage{561}
\mspage{584}

\diagram{
\((62)\qquad\)
\begin{tikzcd}[column sep=small,row sep=large]
 & & L_0\times SN_\rho^*\times SN_\sigma^*
 & & & & L_0 & & & &\\
 & & & & S_0\Gamma_{\rho,\sigma}
 & & & & Z_{\rho,\sigma} & &\\
L_0\times SN_\rho^*\times SN_\sigma^*\times\mathcal G^+
 & & SN_\rho^*\times SN_\sigma^*
 & & & & \mathrm{III}^{\unsure{+|'}}
 & & & &\\
 & & & S_0\mathcal M_{\rho,\sigma}
 & \Gamma_{\rho,\sigma} & & N_\rho^*
 & & & & \{\pm1\}\\
 & SN_\rho^*\times SN_\sigma^*\times\mathcal G^+
 & & & & & N_\sigma^*
 & & \gamma_{\rho,\sigma} & & \widehat{\mathbb Z}^{*}\\
 & & & & & \mathrm{II}
 & \deleted{\unsure{\mathcal G^+|N_\sigma^*}}
 & & & & \{\pm1\}\\
 & & & \mathcal M_{\rho,\sigma}
 & & & & & G_{\rho,\sigma} & &\\
\mathcal G^+ & & & & & & & & & &\\
 & & \mathcal G & & & & & & & &\\
 & & \Vert & & & & & & & &\\
 & & \{1,\tau'\}\mathbin{.}\mathcal G^+
 & & & & & & & &
\arrow[from=1-3,to=1-7,"\mathrm{pr}"]
\arrow[from=1-3,to=2-5,hook]
\arrow[from=1-3,to=3-3,"\mathrm{pr}"]
\arrow[from=1-7,to=2-9,hook]
\arrow[from=2-5,to=2-9,"\theta_Z"]
\arrow[from=2-5,to=4-5]
\arrow[from=2-9,to=5-9,"\delta_Z"]
\arrow[from=3-1,to=1-3,"\mathrm{pr}"]
\arrow[from=3-1,to=5-2,"\mathrm{pr}"']
\arrow[from=3-1,to=4-4,hook]
\arrow[from=3-3,to=4-5,hook]
\arrow[from=4-4,to=2-5]
\arrow[from=4-4,to=7-4]
\arrow[from=4-5,to=4-7,"\theta_\rho"]
\arrow[from=4-5,to=5-7,"\theta_\sigma"']
\arrow[from=4-7,to=5-9,"\delta_\rho"]
\arrow[from=5-7,to=5-9,"\delta_\sigma"]
\arrow[from=5-2,to=3-3,"\mathrm{pr}"']
\arrow[from=5-2,to=7-4,hook]
\arrow[from=5-2,to=6-7,dashed,"\unsure{\mathrm{pr}}"]
\arrow[from=6-7,to=7-9,dashed,hook]
\arrow[from=7-4,to=4-5]
\arrow[from=7-4,to=6-7,dotted]
\arrow[from=7-4,to=7-9,"\theta_G"]
\arrow[from=7-9,to=5-9,"\delta_G"']
\arrow[from=5-9,to=4-11,
 "{\varepsilon_{\gamma,\rho}=\varepsilon'_\gamma}"]
\arrow[from=5-9,to=5-11,"\chi_\gamma"]
\arrow[from=5-9,to=6-11,
 "{\varepsilon_{\gamma,\sigma}=\varepsilon_\gamma}"']
\arrow[from=8-1,to=5-2,hook,"\uline{\mathrm{inj}}"]
\arrow[from=8-1,to=7-4,dashed]
\arrow[from=8-1,to=9-3,hook]
\arrow[from=9-3,to=7-4,hook,"i_{\mathcal G,\mathcal M}"]
\arrow[from=9-3,to=7-9,dashed,hook,"i_{\mathcal G,G}"']
\arrow[from=1-3,to=2-9,phantom,
 "{\deleted{\unsure{\mathrm{IV}}}\ \mathrm{V}}" description,pos=.52]
\arrow[from=4-7,to=5-7,phantom,"\mathrm{I}" description]
\arrow[from=8-1,to=7-4,phantom,"\mathrm{IV}" description,pos=.5]
\end{tikzcd}
}

We have summarized \deleted{\unsure{the main}}\inserted{some of the}
relations between the groups introduced so far in the commutative
diagram (62), which I comment on briefly.

a) All the arrows are either injections \((\hookrightarrow)\)
or surjections \((\longrightarrow)\), with the exception
\inserted{perhaps} of the three arrows issuing from \(\gamma_{\rho,\sigma}\).

b) The ten squares (six faces of a cube, and
\deleted{four}\inserted{\unsure{three|the}} squares marked
\(\mathrm I,\mathrm{II},\mathrm{III}^{\unsure{+|'}},
\mathrm{IV},\mathrm V\)) are cartesian.

c) The \inserted{\unsure{five}} arrows marked \(\mathrm{pr}\)
are projections, canonical from product groups onto partial
products. The arrow marked \(\mathrm{inj}\)
\diagram{
\begin{tikzcd}
\mathcal G^+\arrow[r] &
SN_\rho^*\times SN_\sigma^*\times\mathcal G^+
\end{tikzcd}
}
is the canonical injection.

d) In order not to destroy the commutativity of the diagram,
we have omitted to show in it the arrows
\inserted{(of inclusion, in particular)}
\diagram{
\begin{tikzcd}
Z_{\rho,\sigma}\arrow[r] & G_{\rho,\sigma},
\end{tikzcd}
\qquad
\begin{tikzcd}
N_\rho^*\arrow[r] & G_{\rho,\sigma},
\end{tikzcd}
\qquad
\begin{tikzcd}
N_\sigma^*\arrow[r] & G_{\rho,\sigma}
\end{tikzcd}
}
which appeared in (45).

\mspage{560}

{\small
\begin{center}
560\\
\texttt{https://grothendieck.umontpellier.fr}\\
Université de Montpellier, UMR CNRS 5149, Case courrier 051, Place Eugène Bataillon\\
34095 Montpellier cedex 5 - France
\end{center}
}

\scanpage{562}

e) In order not to overload the diagram, I have not introduced in
the diagram notations for certain useful composite arrows,
such as
\[
(63)\quad
\left\{
\begin{array}{ll}
\delta_\Gamma:\Gamma_{\rho,\sigma}\longrightarrow\gamma_{\rho,\sigma},
&
\deleted{\unsure{\delta_{S_0\Gamma}}:
\unsure{S_0\Gamma_{\rho,\sigma}}\longrightarrow
\unsure{\gamma_{\rho,\sigma}}.}\\
\delta_{S_0\Gamma}:S_0\Gamma_{\rho,\sigma}
\longrightarrow\gamma_{\rho,\sigma}&\\
\delta_{\mathcal M}:\mathcal M_{\rho,\sigma}
\longrightarrow\gamma_{\rho,\sigma}&\\
\delta_{S_0\mathcal M}:S_0\mathcal M_{\rho,\sigma}
\longrightarrow\gamma_{\rho,\sigma}&
\end{array}
\right.
\]
and
\[
(64)\quad
\left\{
\begin{aligned}
\chi_\Gamma&=\chi_\gamma\circ\delta_\Gamma:
\Gamma_{\rho,\sigma}\longrightarrow\widehat{\mathbb Z}^{*}\\
\chi_{S_0\Gamma}&=\chi_\gamma\circ\delta_{S_0\Gamma}:
S_0\Gamma_{\rho,\sigma}\longrightarrow\widehat{\mathbb Z}^{*}\\
\chi_{\mathcal M}&=\chi_\gamma\circ\delta_{\mathcal M}:
\mathcal M_{\rho,\sigma}\longrightarrow\widehat{\mathbb Z}^{*}\\
\chi_{S_0\mathcal M}&=\chi_\gamma\circ\delta_{S_0\mathcal M}:
S_0\mathcal M_{\rho,\sigma}\longrightarrow\widehat{\mathbb Z}^{*}
\end{aligned}
\right.
\]
and the composites
\(\varepsilon_{?,\rho},\varepsilon_{?,\sigma}:?\longrightarrow\{\pm1\}\),
where \(?\) can take \deleted{\unsure{any}} value
\(Z,G\) (already introduced) \(N_\rho^*,N_\sigma^*\)
(we will have
\(\varepsilon_{N_\rho,\rho}=\varepsilon_\rho
\deleted{\text{\unsure{simply}}},\quad
\varepsilon_{N_\sigma,\sigma}=\varepsilon_\sigma
\deleted{\text{\unsure{simply}}}\)),
\(\Gamma,S_0\Gamma,\mathcal M,S_0\mathcal M\).
The kernel of \(\chi_?\) is denoted \(?^+\), that of
\(\varepsilon_{?,\rho}\) is denoted \(?'\), that of
\(\varepsilon_{?,\sigma}\) is denoted \(?''\),
we write
\[
?^\circ\doteq?'\cap?'',\qquad
?^{+'}=?^+\cap?',\quad
?^{+''}=?^+\cap?'',
\]
\[
?^{+,\circ}=?^+\cap?^\circ,\quad
S?=\operatorname{Ker}\delta_?,
\]
\deleted{and} hence \(\deleted{\unsure{S?\subset}}\)

\margin{\deleted{there is a small incompatibility with the notation
\(N'_\rho,N'_\sigma\) already used in another \unsure{sense}.}}

\[
(65)
\]
\diagram{
\begin{tikzcd}[row sep=1.15em,column sep=1.6em]
 & & ?^{'+} & &\\
\deleted{\unsure{S?\subset}}\quad ?^{\circ+}
 & & & & ?^+\ \deleted{\text{\unsure{we}}}\\
 & & ?^{''+} & &\\
 & & \unsure{?'} & &\\
?^\circ & & & & ?\\
 & & ?'' & &
\arrow[from=2-1,to=1-3,hook]
\arrow[from=1-3,to=2-5,hook]
\arrow[from=2-1,to=3-3,hook]
\arrow[from=3-3,to=2-5,hook]
\arrow[from=2-1,to=5-1,hook]
\arrow[from=1-3,to=4-3,hook,shift left=.7ex]
\arrow[from=3-3,to=6-3,hook,shift right=.7ex]
\arrow[from=2-5,to=5-5,hook]
\arrow[from=5-1,to=4-3,hook]
\arrow[from=5-1,to=6-3,hook]
\arrow[from=4-3,to=5-5,hook]
\arrow[from=6-3,to=5-5,hook]
\end{tikzcd}
}
``cartesian'' cube.

f) Let us also recall that we have commutativity in
\[
(66)
\]
\diagram{
\begin{tikzcd}[row sep=.8em,column sep=2em]
 & \{\pm1\}\arrow[dr] &\\
\mathcal G
\arrow[ur,"{\varepsilon_{\mathcal G,\rho}}"]
\arrow[rr,"\chi_{\mathcal G}"]
\arrow[dr,"{\varepsilon_{\mathcal G,\sigma}}"']
 & & \widehat{\mathbb Z}^{*}\\
 & \{\pm1\}\arrow[ur] &
\end{tikzcd}
}
, i.e. \(\varepsilon_{\mathcal G}:\mathcal G\longrightarrow\{\pm1\}\),

\mspage{561}

{\small
\begin{center}
561\\
\texttt{https://grothendieck.umontpellier.fr}\\
Université de Montpellier, UMR CNRS 5149, Case courrier 051, Place Eugène Bataillon\\
34095 Montpellier cedex 5 - France
\end{center}
}